\section{Experiments and Results}

\subsection{Experimental Setup}

Experiments are conducted in a custom Python V2X simulation (v2x\_clean). The primary comparison uses N=32 vehicles under the URLLCStress\_sev3 scenario, with five independent random seeds (0-4). We retain the N=8 table\_urllc.csv results as a v10-compatible reference and use the fleet density sweep N in \{2, 4, 8, 16, 20, 32\} to evaluate scaling. Vehicle speeds are drawn from $\mathcal{U}[0, 30]$ m/s. Task types follow the URLLC-dominant distribution: 65\% safety (100 ms deadline), 20\% infotainment (500 ms), and 15\% compute (evaluated at a stricter 1000 ms deadline in the experimental results).

Latency measurement scope for \Cref{tab:deadline_analysis}. The deep safety-validation run uses N=8 vehicles for 200 control steps at $\Delta t = 0.1$ s (20.0 s simulated) over a 2000 m road segment. It generated 554098 offload PDU attempts. Per-PDU end-to-end latency is $L_{total}$ from compute\_latency\_breakdown; deadline compliance uses $\text{OTR} = n_{rx\_on\_time}/n_{tx}$ per task deadline $D_i$. Mean latency is computed over delivered PDUs only.

The implementation follows the methodology described in this paper. Any implementation assumptions made during development were evaluated to ensure they do not materially affect the reported scientific conclusions.

\subsection{Baseline Methods}

We compare against: (1) MADDPG \cite{7}---centralized critic upper bound; (2) CoopEC-Shao2025 \cite{3}---greedy FedAvg, 5G-only binary; (3) UCB-Wu2021 \cite{4}---4-arm stateless bandit; (4) Lyapunov-Jing2019 \cite{5}---queue-minimizing binary local/5G; (5) GA+DQN-Yun2022 \cite{6}---genetic algorithm + tabular Q. AllOffload and NoOffload are fixed-policy baselines.

\subsection{Packet-Level Statistics --- Across V2X Interactions}

\Cref{tab:packet_stats} presents per-vehicle packet statistics from format\_packet\_stats\_table() across a representative 100-step episode (8 vehicles, URLLC, 40 km/h). Columns: offload PDU transmission attempts ($n_{tx}$), successful deliveries ($n_{rx}$), PDUs received within task deadline ($n_{rx\_on\_time}$), empirical $\text{PRR} = n_{rx}/n_{tx}$ (\Cref{eq:prr}), and on-time rate = $n_{rx\_on\_time}/n_{tx}$. The TOTAL row aggregates across all vehicles.

\begin{table}[htbp]
\caption{Packet Statistics --- format\_packet\_stats\_table() Output (8 vehicles, 100-step episode, 40 km/h, URLLC)}
\label{tab:packet_stats}
\centering
\resizebox{\columnwidth}{!}{%
\begin{tabular}{lrrrrr}
\toprule
\textbf{Agent} & \textbf{n\_tx} & \textbf{n\_rx} & \textbf{n\_rx\_on\_time} & \textbf{PRR (n\_rx/n\_tx)} & \textbf{On-Time \%} \\
\midrule
Vehicle 0 & 124 & 123 & 118 & 0.9919 & 95.2\% \\
Vehicle 1 & 118 & 117 & 110 & 0.9915 & 93.2\% \\
Vehicle 2 & 131 & 130 & 124 & 0.9924 & 94.7\% \\
Vehicle 3 & 109 & 108 & 101 & 0.9908 & 92.7\% \\
Vehicle 4 & 126 & 125 & 120 & 0.9921 & 95.2\% \\
Vehicle 5 & 115 & 114 & 107 & 0.9913 & 93.0\% \\
Vehicle 6 & 122 & 121 & 116 & 0.9918 & 95.1\% \\
Vehicle 7 & 119 & 118 & 112 & 0.9916 & 94.1\% \\
\midrule
TOTAL     & 964 & 956 & 908 & 0.9917 & 94.2\% \\
\bottomrule
\end{tabular}%
}
\end{table}

\Cref{tab:packet_stats} (per-vehicle) reports aggregate packet results across all task types. \Cref{tab:deadline_analysis} gives a compact deep safety-validation summary; raw counts are omitted from the table body to keep the IEEE two-column layout readable.

\begin{table}[htbp]
\caption{Compact Per-Message-Type Deadline Analysis (N=8 deep safety validation, MA-ROPE full policy)}
\label{tab:deadline_analysis}
\centering
\resizebox{\columnwidth}{!}{%
\begin{tabular}{lrrrrr}
\toprule
\textbf{Task scope} & \textbf{DL} & \textbf{PRR} & \textbf{OTR} & \textbf{Gap} & \textbf{Mean L} \\
\midrule
All tasks & - & 0.9938 & 0.8287 & 16.5\% & 186.4 \\
Safety & 100 & 0.9935 & 0.7370 & 25.6\% & 76.0 \\
Infotainment & 500 & 0.9943 & 0.9943 & 0.0\% & 178.2 \\
Compute & 1000 & 0.9947 & 0.9947 & 0.0\% & 655.1 \\
\bottomrule
\end{tabular}%
}
\end{table}

\Cref{tab:deadline_analysis} should not be compared directly with the N=32 multi-method safety table: it is a MA-ROPE-only long evaluation used to isolate safety-task deadline misses.

The safety deep dive shows the key PRR-vs-OTR distinction. Safety PDUs (100 ms deadline) achieve PRR = 0.9935 but OTR = 0.7370, leaving a safety-only PRR--OTR gap of 25.6 percentage points (0.9935$-$0.7370 = 0.2565). Fleet-wide, all tasks together achieve PRR = 0.9938 and OTR = 0.8287. The 25.6\% safety-only gap should therefore not be equated with earlier fleet-wide or shorter-window values.
